%%%%%%%%%%  scan idx 458  |  folio 445  %%%%%%%%%%
\[
\begin{tikzcd}[column sep=large,row sep=large]
Y' \arrow[r,"j"] \arrow[d,"g"'] & X' \arrow[d,"f"]\\
Y \arrow[r,"i"] & X.
\end{tikzcd}
\]
On
suppose que \(X\) est quasi-compact et a un faisceau ample. Alors, on a la suite exacte:
\[
0\longrightarrow K^\circ(Y)\xrightarrow{u}K^\circ(Y')\times K^\circ(X)
\xrightarrow{v}K^\circ(X'),
\]
où on pose, pour \(y\in K^\circ(Y)\):
\[
u(y)=\bigl(g^*(y).\lambda_{-1}(F),-i_*(y)\bigr),
\]
et pour \(y'\in K^\circ(Y')\), \(x\in K^\circ(X)\):
\[
v(y',x)=j_*(y')+f^*(x),
\]
les notations étant celles de 3.1 et 3.3. En outre, \(u\) admet un inverse à gauche \(u'\), défini par \(u'(y',x)=g_*(y')\).

Si de plus \(X\) est noethérien, et \(Y\) régulier, alors \(v\) est surjectif, donc on a la suite exacte:
\[
0\longrightarrow K^\circ(Y)\xrightarrow{u}K^\circ(Y')\times K^\circ(X)
\xrightarrow{v}K^\circ(X')\longrightarrow0.
\]

Tout d'abord, \(u\) admet l'inverse à gauche \(u'\), et a fortiori est injectif. En effet, on a:
\[
g_*\bigl(g^*(y).\lambda_{-1}(F)\bigr)
=y.g_*(\lambda_{-1}(F))
\]
d'après la formule de projection. Par ailleurs, \(g_*(\lambda_{-1}(F))=1\) (VI 5.9), d'où le résultat.

Ensuite, \(v\circ u=0\). Si \(y\in K^\circ(Y)\), on a en effet:
\[
v\circ u(y)
=v\bigl(g^*(y)\lambda_{-1}(F),-i_*(y)\bigr)
=j_*\bigl(g^*(y)\lambda_{-1}(F)\bigr)-f^*(i_*(y))=0,
\]
d'après 3.4.

Il reste donc à montrer que \(\Ker(v)\subset \Im(u)\). Soient donc \(x\in K^\circ(X)\), \(y'\in K^\circ(Y')\), tels que:
